\chapter{Ankle Pain}

\section*{Definition}
Pain or discomfort localized to the ankle joint and surrounding structures, including the medial and lateral malleoli, talocrural joint, and supporting ligaments and tendons.

\section*{When to See a Doctor}

\begin{itemize}
 \item Inability to take four steps after injury, severe swelling or deformity, numbness or a cold foot, fever with ankle pain, or pain after significant trauma
 \item Persistent pain, swelling, or instability that is not improving over several weeks
\end{itemize}

\section*{How It Develops}
Ankle pain most commonly follows a ligament injury, often an inversion sprain affecting the lateral ligament complex. It can also arise from bone, tendon, cartilage, joint, nerve, or skin disease. Persistent symptoms may reflect instability, tendinopathy, post-traumatic arthritis, or an osteochondral injury.

\section*{Causes}
\begin{itemize}
 \item Acute ankle sprain (lateral, medial, or syndesmotic)
 \item Ankle fracture (malleolar, bimalleolar, trimalleolar)
 \item Achilles tendinopathy or rupture
 \item Peroneal tendinopathy or subluxation
 \item Posterior tibial tendinopathy
 \item Osteochondritis dissecans of the talus
 \item Ankle osteoarthritis (post-traumatic)
 \item Gout or pseudogout
 \item Cellulitis or septic arthritis
 \item Tarsal coalition (congenital, presenting in adolescence)
\end{itemize}

\section*{Related Symptoms}
\begin{itemize}
 \item Foot pain
 \item Leg pain
 \item Swelling
 \item Bruising
 \item Instability
\end{itemize}


\section*{Medical Evaluation}
\begin{itemize}
 \item The clinician will assess ability to bear weight and tenderness at the posterior edges or tips of both malleoli, the base of the fifth metatarsal, and the navicular, applying the Ottawa Ankle and Foot Rules when appropriate.
 \item X-rays are obtained when fracture criteria or another indication is present; weight-bearing views are used selectively. Further imaging is not needed for every sprain.
 \item MRI is indicated when osteochondral defects of the talus, occult fractures, or ligamentous injuries are suspected.
 \item Referral to an orthopaedic surgeon or sports medicine specialist is arranged for severe sprains, fractures, or chronic instability.
\end{itemize}

\section*{Treatment Options}
\begin{itemize}
 \item Functional rehabilitation with progressive range-of-motion, strengthening (peroneals, calf raises), and proprioceptive training is the mainstay of ankle sprain management.
 \item Immobilisation with a walking boot or cast is used for stable fractures, followed by graduated weight-bearing and physiotherapy.
 \item Surgical repair is indicated for unstable fractures, syndesmotic injuries (high ankle sprains), or chronic lateral ankle instability failing conservative therapy.
 \item Corticosteroid injections are not routinely used for Achilles tendinopathy and can increase rupture risk; any injection near a tendon requires specialist risk assessment.
\end{itemize}

\section*{Self-Care and Home Management}
\begin{itemize}
 \item Protect the ankle, use ice wrapped in cloth for up to 20 minutes at a time if helpful, consider comfortable compression, and elevate it during the first days. Avoid prolonged complete rest.
 \item Take over-the-counter NSAIDs (ibuprofen, naproxen) for pain and inflammation unless contraindicated.
 \item Use an ankle brace or supportive wrap during weight-bearing activities until pain subsides.
 \item Begin gentle range-of-motion and progressive weight-bearing as tolerated, then strengthening and balance exercises; a clinician or physiotherapist can guide timing for more severe injuries.
\end{itemize}

\section*{References}
\begin{thebibliography}{99}
\bibitem{ankle_pain:ref1} Wolfe MW, Uhl TL, et al. ``Management of ankle sprains.'' \textit{Am Fam Physician}. 2001;63(1):93-104.
\bibitem{ankle_pain:ref2} Doherty C, Bleakley C, et al. ``Treatment and prevention of acute and recurrent ankle sprain.'' \textit{Sports Med}. 2017;47(1):123-144.
\bibitem{ankle_pain:ref3} Stiell IG, Greenberg GH, et al. ``The Ottawa Ankle Rules: a clinical decision rule for radiography of ankle injuries.'' \textit{JAMA}. 1994;271(11):827-832.
\bibitem{ankle_pain:ref4} DeLee JC, Drez D, et al. \textit{DeLee \& Drez\'s Orthopaedic Sports Medicine}, 5th ed. Elsevier, 2022.
\bibitem{ankle_pain:ref5} UpToDate. ``Ankle sprain in adults: evaluation and management.'' Wolters Kluwer, 2023.
\bibitem{ankle_pain:ref6} Martin RL, et al. Ankle stability and movement coordination impairments: lateral ankle ligament sprains revision. \textit{J Orthop Sports Phys Ther}. 2021;51(4):CPG1--CPG80.
\end{thebibliography}
